\chapter{From EKF-SLAM to the MSCKF}
\label{ch:msckfintro}

\begin{objectives}
\begin{itemize}
  \item Understand how EKF-SLAM uses visual landmarks, and why it scales poorly.
  \item Understand the central idea of the MSCKF: constraints between multiple past poses instead of landmarks.
  \item See why past poses must be kept in the state, with their correlations.
  \item Relate the MSCKF to sliding-window optimization.
\end{itemize}
\end{objectives}

% ------------------------------------------------------------------------------
\section{EKF-SLAM}

\index{EKF-SLAM}
The most direct way to use a camera in an EKF is to put the landmarks in the state. With $L$ landmarks,
\begin{equation}
  x = \big(x_I,\ p_{f_1},\ \dots,\ p_{f_L}\big), \qquad \dim\delta x = 15 + 3L,
\end{equation}
and each image observation $z = \pi(p_c(x_I, p_{f_j})) + n$ is an ordinary EKF measurement of the IMU state and one
landmark. This is EKF-SLAM, and MonoSLAM~\cite{davison2007monoslam} showed it can run in real time with a single camera.

\subsection{How it works}

\paragraph{Prediction.} Landmarks are static, so only the IMU block is propagated: exactly the block structure of clones in \cref{ch:errordyn},
$P_{II} \gets \Phi P_{II}\Phi\T + Q_d$, $P_{IL} \gets \Phi P_{IL}$, with the landmark block unchanged.

\paragraph{Update.} An observation of landmark $j$ has a Jacobian with two non-zero blocks: one for the IMU pose and one for landmark $j$,
$H = [H_I\ \ 0\ \cdots\ 0\ \ H_{f_j}\ \ 0\ \cdots]$. The gain $K = PH\T S\inv$ is nevertheless dense: every landmark correlated with the robot or with landmark $j$ is corrected.
This is how EKF-SLAM ``closes loops'': re-observing an old landmark corrects the robot, which through the correlations corrects every other landmark. It is also
what makes it expensive.

\paragraph{Landmark initialization.} A new landmark enters the state as a function of the current state and the first measurement, $p_f = g(x_I, z)$. This is the
state augmentation of \cref{prop:augmentation}: $P \gets JPJ\T + J_zRJ_z\T$ with $J = [I;\ \partial g/\partial x]$ and $J_z = [0;\ \partial g/\partial z]$. For a monocular camera, $g$
needs a depth that a single measurement does not provide, which leads to the initialization problem below.

\subsection{Why it scales poorly}

EKF-SLAM has three structural problems.
\begin{enumerate}
  \item \textbf{Cost.} The covariance has $(15 + 3L)^2$ entries, and every update touches all of them: the Joseph update costs $\mathcal{O}(n^2)$ per scalar measurement at
        best. Tracking 200 landmarks gives a $615\times615$ covariance; updating it at 20~Hz with hundreds of observations per frame is expensive, and the cost
        grows quadratically with the map.
  \item \textbf{Initialization.} A single monocular observation says nothing about depth. A new landmark must be
        initialized with a huge depth uncertainty, which is highly non-Gaussian in Euclidean coordinates, or its
        initialization must be delayed until it has been seen from several viewpoints. Inverse-depth
        parameterizations~\cite{civera2008inverse} help, but add complexity.
  \item \textbf{Housekeeping.} Landmarks that leave the field of view must either stay in the state, growing it, or be
        marginalized, which in covariance form is cheap but discards their future value.
\end{enumerate}

% ------------------------------------------------------------------------------
\section{The multi-state constraint}
\index{multi-state constraint}\index{MSCKF}

Consider a feature observed in $M$ images taken at times $t_1, \dots, t_M$. Each observation constrains the pose at its
time and the feature position. Taken together, after the unknown feature position is eliminated, the $2M$ observations
impose $2M - 3$ constraints \emph{on the $M$ poses alone}. This is the \emph{multi-state constraint}: geometric information
about how the camera moved between the images, with the scene point serving only as an intermediary.

\begin{keyidea}[the MSCKF in one sentence]
Keep a sliding window of past poses in the state instead of landmarks; when a feature track is complete,
triangulate it, linearize its observations, and eliminate the feature by projecting onto the left nullspace of its
Jacobian, leaving a residual that depends only on the poses; then update.
\end{keyidea}

Mourikis and Roumeliotis~\cite{mourikis2007msckf} showed that this uses all the information that the observations carry about
the poses (under linearization), at a cost linear in the number of features. Feature positions are never part of the
state, so there is no landmark initialization problem inside the filter, and features leaving the view simply trigger an
update.

\section{Why past poses must be in the state}

Why not keep only the current IMU state and use the constraint between the current pose and the past ones? Because the
past poses are \emph{estimated} quantities, with errors correlated with the current state's error: they were all produced by
the same propagation from the same measurements. A measurement relating the current pose to a past pose estimate, which
treats that past estimate as fixed or independent, double counts information and makes the filter overconfident.

The cure is \emph{stochastic cloning}~\cite{roumeliotis2002cloning}. When an image is captured, a copy of the current pose (a
clone) is appended to the state, together with its exact correlations to everything else. From then on the clone is a
state variable whose estimate can be corrected by later measurements, and whose correlation with the evolving IMU state is
propagated exactly. \Cref{ch:cloning} develops this in detail.

\begin{figure}[t]
\centering
\begin{tikzpicture}[scale=1.0]
  % trajectory
  \draw[->, thick, black!50] (-0.5,-0.2) .. controls (2,0.6) and (5,0.7) .. (9.4,0.1) node[right, font=\footnotesize] {time};
  % clones
  \foreach \x/\y/\a/\n in {0.3/0.02/12/1, 1.8/0.33/9/2, 3.3/0.5/5/3, 4.8/0.56/0/4, 6.3/0.5/-5/5, 7.8/0.33/-10/6} {
    \begin{scope}[shift={(\x,\y)}, rotate=\a]
      \draw[fill=primarylight, draw=primary] (-0.2,0) -- (0.2,0) -- (0.3,0.42) -- (-0.3,0.42) -- cycle;
    \end{scope}
    \node[font=\scriptsize, below=2pt] at (\x,\y) {$C_\n$};
  }
  % current IMU
  \fill[conv] (9.0,0.15) circle (2.5pt) node[below=3pt, font=\scriptsize] {IMU state};
  % features
  \coordinate (F1) at (2.4,3.2); \coordinate (F2) at (5.6,3.6); \coordinate (F3) at (7.6,2.9);
  \fill[plotgreen] (F1) circle (2.3pt) node[above] {\footnotesize $f_1$};
  \fill[plotgreen] (F2) circle (2.3pt) node[above] {\footnotesize $f_2$};
  \fill[plotgreen] (F3) circle (2.3pt) node[above] {\footnotesize $f_3$};
  \foreach \x/\y in {0.3/0.44, 1.8/0.75, 3.3/0.92} \draw[plotgreen, dashed] (\x,\y) -- (F1);
  \foreach \x/\y in {3.3/0.92, 4.8/0.98, 6.3/0.92, 7.8/0.75} \draw[plotgreen!60!black, dashed] (\x,\y) -- (F2);
  \foreach \x/\y in {6.3/0.92, 7.8/0.75} \draw[plotorange, dashed] (\x,\y) -- (F3);
  \draw[decorate, decoration={brace, amplitude=5pt, mirror}] (0,-0.55) -- (8.1,-0.55) node[midway, below=6pt, font=\footnotesize] {sliding window of $N$ clones (poses at image times), all in the state};
\end{tikzpicture}
\caption{The MSCKF state: the current IMU state plus clones of past camera poses. Feature $f_1$ is seen by clones 1--3
and constrains them; $f_2$ by clones 3--6. Neither feature is in the state. Feature $f_3$ has only two observations so far
and keeps being tracked.}
\label{fig:msckf-window}
\end{figure}

% ------------------------------------------------------------------------------
\section{The MSCKF state}

The project's state (AGENTS.md C3) is the IMU state plus $N \le$ \code{kMaxClones} clones, each holding the IMU pose
$(p_{wb}, q_{wb})$ at an image timestamp:
\begin{equation}
  x = \big(\,p,\ v,\ q,\ b_a,\ b_g,\ \ p_1, q_1,\ \dots,\ p_N, q_N\,\big),\qquad \dim\delta x = 15 + 6N .
\end{equation}
The camera pose of clone $i$ is derived through the fixed extrinsics,
$R_{wc_i} = R_{wb_i}R_{bc}$ and $p_{wc_i} = p_{wb_i} + R_{wb_i}p_{bc}$. With \code{kMaxClones} $= 10$, the error state has 75
dimensions, independent of how many features are tracked.

% ------------------------------------------------------------------------------
\section{Relation to sliding-window optimization}

Optimization-based VIO (OKVIS~\cite{leutenegger2015okvis}, VINS-Mono~\cite{qin2018vins}) keeps a similar window of poses and
minimizes the sum of squared reprojection and IMU residuals over poses \emph{and} landmarks, relinearizing at every
iteration, then marginalizes old states into a Gaussian prior. The MSCKF update is equivalent to a single Gauss--Newton
iteration of such a problem, with the landmarks eliminated by a Schur complement and the prior given by the filter's
covariance. Eliminating a landmark by nullspace projection and by Schur complement give the same information about the
poses (\cref{thm:nullspace-optimal}). The differences are that the MSCKF never relinearizes past measurements, and
processes each feature track exactly once, which makes it cheaper and more predictable, at some cost in accuracy under
strong nonlinearity.

\section{Variants}

The basic MSCKF has been refined in many directions. MSCKF 2.0~\cite{li2013msckf2} added First-Estimates Jacobians for
consistency and online calibration of extrinsics and time offset. Hybrid filters keep long-lived features in the state as
SLAM landmarks while processing short tracks MSCKF-style. MSCKF-VIO~\cite{sun2018msckfvio} uses a stereo camera and
two-point RANSAC for fast flight, and OpenVINS~\cite{geneva2020openvins} is a modular research platform that implements most of
these options and documents them carefully.

\begin{takeaways}
\begin{itemize}
  \item EKF-SLAM keeps landmarks in the state; its cost grows quadratically with the map.
  \item The MSCKF keeps past poses (clones) instead; each feature track is turned into constraints on those poses.
  \item Past poses must be in the state, with their correlations, or information is double counted.
  \item The MSCKF update is equivalent to one Gauss--Newton step of sliding-window bundle adjustment with landmarks
        marginalized.
\end{itemize}
\end{takeaways}

\begin{exercises}
  \item Estimate the memory and the per-update flop count of EKF-SLAM with 200 landmarks and of the MSCKF with 10 clones.
  \item A feature is observed in $M$ images. How many scalar constraints does it impose on the poses? Why does the count
        start being useful only at $M = 2$, and why does the project require $M \ge 3$?
  \item Explain in your own words why using a past pose estimate as if it were known exactly makes the filter
        overconfident. Construct a scalar example.
\end{exercises}
